\documentclass[../DoAn.tex]{subfiles}
\begin{document}

Chương \ref{chapter:Experiment} đã trình bày thiết kế và kết quả xây dựng hệ thống, trong đó một số giải pháp kỹ thuật chỉ được nhắc tới ngắn gọn. Chương này trình bày chi tiết bảy giải pháp mà đồ án đóng góp, mỗi giải pháp gồm vấn đề đặt ra, cách giải quyết và kết quả đo được: tìm kiếm lai cho văn bản pháp lý tiếng Việt (phần \ref{section:5.1}), tác tử tự kiểm tra quy định hiện hành (phần \ref{section:5.2}), đồ thị pháp luật cho dẫn chiếu và kiểm tra ngoại lệ (phần \ref{section:5.graph}), kiểm chứng căn cứ trong rà soát văn bản (phần \ref{section:5.3}), bộ nhớ dài hạn bảo vệ quyền riêng tư (phần \ref{section:5.4}), quy trình cập nhật văn bản và theo dõi hiệu lực (phần \ref{section:5.5}) và trực quan hoá hoạt động của tác tử (phần \ref{section:5.6}).

\section{Tìm kiếm lai cho văn bản pháp lý tiếng Việt}
\label{section:5.1}

\subsection{Vấn đề}
Phiên bản đầu tiên của hệ thống chỉ dùng tìm kiếm vector trên tên điều luật: lấy 30 điều có tên gần nhất với câu hỏi, gán cùng một điểm cơ sở 0,8 cho mọi kết quả, rồi cộng thêm 0,05 điểm cho mỗi từ của câu hỏi xuất hiện trong điều luật. Cách làm này có ba nhược điểm. Thứ nhất, tên điều luật rất ngắn và mang tính khái quát (``Quyền và nghĩa vụ của người lao động''), nên câu hỏi chứa thuật ngữ nằm trong nội dung như ``thai sản'', ``án treo'', ``nồng độ cồn'' dễ bị bỏ lỡ. Thứ hai, việc gán cùng một điểm cơ sở làm mất thông tin thứ hạng của Milvus, thứ tự cuối cùng chỉ còn phụ thuộc vào việc đếm từ. Thứ ba, phép đếm từ tính cả các hư từ như ``của'', ``là'', ``được'', khiến điều luật dài và nhiều hư từ được ưu tiên.

Không thể giải quyết vấn đề bằng chức năng tìm kiếm toàn văn sẵn có của MySQL, vì chức năng này mặc định bỏ qua các từ dưới ba ký tự, tức là bỏ phần lớn âm tiết tiếng Việt như ``án'', ``tù'', ``xe''. Milvus phiên bản 2.4 cũng chưa có BM25 tích hợp.

\subsection{Giải pháp}
Đồ án xây dựng một bộ tìm kiếm lai gồm ba bước.

\textbf{Bước 1: BM25 với bộ tách từ cho tiếng Việt.} Bộ tách từ chuyển văn bản về chữ thường, chuẩn hoá Unicode dạng NFC rồi tách theo âm tiết. Các âm tiết được chia thành hai nhóm từ dừng. Nhóm hư từ (``và'', ``của'', ``là'', ``tôi'', ``mình'', ``không''\ldots) bị loại khỏi cả âm tiết đơn lẫn cặp âm tiết. Nhóm từ dừng mềm (``điều'', ``khoản'', ``luật'', ``quy'', ``làm''\ldots) bị loại khi đứng một mình nhưng được giữ trong cặp âm tiết, vì chúng mang nghĩa trong từ ghép như ``điều\_kiện'', ``quy\_định'', ``làm\_việc''. Bên cạnh âm tiết đơn, bộ tách từ sinh thêm mọi cặp âm tiết liền kề không chứa hư từ; do phần lớn từ ghép tiếng Việt có hai âm tiết, cách này nắm được ``thai\_sản'', ``lao\_động'' mà không cần bộ tách từ dựa trên từ điển. Tên điều được lặp ba lần để tăng trọng số. Mỗi tài liệu dài tối đa 6.000 ký tự; điều dài hơn được tách theo khoản thành nhiều tài liệu và khi tìm kiếm thì gộp về điều với điểm cao nhất, nhờ vậy các khoản cuối của điều dài vẫn tìm được mà điều dài không chiếm ưu thế về tần suất từ. Chỉ mục gồm 68.086 tài liệu từ 66.706 điều, được xây dựng một lần, lưu đệm xuống đĩa và nạp vào bộ nhớ khi khởi động.

\textbf{Bước 2: Hợp nhất thứ hạng.} Kết quả của tìm kiếm vector trên tên điều và BM25 (mỗi nguồn 50 kết quả) được hợp nhất bằng RRF theo công thức \eqref{eq:rrf} với $k = 60$, trọng số $w_{\mathrm{BM25}} = 1{,}0$ và $w_{v} = 0{,}5$. Điểm RRF được chuẩn hoá về khoảng $(0, 1]$ bằng cách chia cho điểm lớn nhất có thể đạt được, tức điểm của một điều đứng đầu ở mọi nguồn có kết quả:
\begin{equation}\label{eq:rrfnorm}
    s_{\mathrm{base}}(d) = \frac{\mathrm{RRF}(d)}{\sum_{s \in S'} w_s / (k + 1)}
\end{equation}
với $S'$ là tập nguồn trả về ít nhất một kết quả. Điểm chuẩn hoá cho các bước sau một thang đo chung, kể cả khi một nguồn không có kết quả. Trọng số $w_{v} = 0{,}5$ được chọn qua thực nghiệm dò trọng số trình bày ở phần kết quả.

\textbf{Bước 3: Xếp hạng lại theo độ phủ từ có nghĩa.} Điểm cuối cùng của mỗi điều được tính như sau:
\begin{equation}\label{eq:rerank}
    s(d) = (1-\lambda)\cdot s_{\mathrm{base}}(d) + \lambda\cdot \frac{\sum_{t \in T_q} c(t, d)}{2\,|T_q|}
\end{equation}
trong đó $s_{\mathrm{base}}$ là điểm RRF đã chuẩn hoá, $T_q$ là tập âm tiết và cặp âm tiết có nghĩa của câu hỏi, $c(t,d)$ bằng 2 nếu $t$ có trong tên điều, bằng 1 nếu chỉ có trong nội dung và bằng 0 nếu không có, và $\lambda = 0{,}3$. Khác với cách đếm cũ, độ phủ bị chặn trong khoảng $[0, 1]$ và không tính hư từ, nên không thể lấn át điểm truy xuất.

Ngoài cấu hình trên, đồ án thử nghiệm thêm bốn biến thể: tìm kiếm vector trên 147.641 đoạn nội dung kèm tiêu đề (phần \ref{section:4.data}), dùng riêng hoặc làm nguồn thứ ba của RRF; cộng thêm điểm cho các điều nối với kết quả đầu qua đồ thị; và kèm các điều được dẫn chiếu vào cuối danh sách kết quả.

\subsection{Kết quả đạt được}
Các cấu hình được đánh giá bằng hai phép đo bổ sung cho nhau (phần \ref{section:3.5}). Phép đo thứ nhất dùng 50 câu hỏi thực tế, lấy từ lịch sử hội thoại của hệ thống và bổ sung một số câu viết tay; 10 kết quả đầu của mọi cấu hình được gộp lại và chấm mù bằng LLM (Claude Sonnet) theo thang 0, 1, 2. Phép đo thứ hai là đánh giá theo mục đã biết trên 100 câu hỏi do GPT-4o-mini sinh từ nội dung của 100 điều chọn ngẫu nhiên, đáp án là chính điều đó. Vì câu hỏi được sinh từ nội dung điều nên thường lặp lại từ ngữ của điều, phép đo này thiên về BM25 và cần được đọc cùng phép đo thứ nhất. Thiết lập chi tiết và thử nghiệm ban đầu trên năm câu hỏi được trình bày ở Phụ lục \ref{appendix:A}. Bảng \ref{tab:kq_truyxuat} trình bày kết quả.

\input{tables/tab_truyxuat}

So với cấu hình ban đầu, cấu hình cuối nâng P@10 từ 0,368 lên 0,544, MRR từ 0,280 lên 0,471, tỷ lệ kết quả đầu tiên liên quan từ 52\% lên 70\% và Hit@10 theo mục đã biết từ 0,32 lên 0,69. So từng câu hỏi theo nDCG@10, cấu hình cuối tốt hơn cấu hình ban đầu ở 35 câu, bằng ở 6 câu và kém ở 9 câu; kiểm định dấu hai phía cho $p < 0{,}001$ (Bảng \ref{tab:pl_sign}). Vector đoạn nội dung khi dùng riêng chỉ nhỉnh hơn vector tên điều và kém xa BM25, và khi thêm vào làm nguồn thứ ba thì thắng 25, thua 23 câu so với cấu hình hai nguồn ($p = 0{,}885$, Bảng \ref{tab:pl_sign_hybrid}), nên không được dùng. Cộng điểm đồ thị và kèm điều được dẫn chiếu cũng không làm thay đổi đáng kể các chỉ số xếp hạng; giá trị của đồ thị nằm ở việc bổ sung ngữ cảnh cho tác tử, được trình bày ở phần \ref{section:5.graph}.

Bảng \ref{tab:trongso} trình bày thực nghiệm dò trọng số vector, trong đó mọi dòng thuộc cùng một lần chạy nên nDCG@10 so được với nhau.

\input{tables/tab_trongso}

Giảm $w_{v}$ làm P@10 và Hit@10 theo mục đã biết tăng dần, đồng thời MRR giảm khi $w_{v}$ xuống dưới 0,5. Đồ án chọn $w_{v} = 0{,}5$ vì có MRR cao nhất và không kém $w_{v} = 1{,}0$ ở chỉ số nào. Tuy nhiên, mọi khác biệt giữa các mức trọng số, cũng như giữa tìm kiếm lai và BM25 đơn lẻ (23 thắng, 23 thua), đều không có ý nghĩa thống kê với 50 câu hỏi; BM25 đơn lẻ có P@5, P@10, Top1 và các chỉ số theo mục đã biết cao hơn, còn tìm kiếm lai có MRR cao hơn. Kết quả chưa đủ để khẳng định tìm kiếm lai tốt hơn BM25 trên dữ liệu này.

Hit@1 theo mục đã biết của tìm kiếm lai (0,27) thấp hơn BM25 (0,42) phản ánh hệ quả của RRF đã nêu ở phần \ref{section:3.3}. Ví dụ, với câu hỏi ``Lao động nữ đang mang thai có bị sa thải không?'', Điều 137 Bộ luật Lao động 2019 (Bảo vệ thai sản) đứng hạng 1 ở BM25 nhưng hạng 97 ở vector tên điều, ngoài 50 kết quả được hợp nhất. Điều này chỉ nhận điểm RRF từ BM25 là $1/61 \approx 0{,}0164$, trong khi một điều đứng hạng $r$ ở cả hai nguồn nhận $1{,}5/(60 + r)$, lớn hơn khi $r \le 31$. Kết quả là Điều 137 xuống hạng 16 sau RRF và hạng 13 sau xếp hạng lại, nên không có trong 10 kết quả trả cho tác tử; với $w_{v} = 1{,}0$, điều này ở hạng 17 và 19 (số liệu trong \texttt{evaluation/\allowbreak{}results/\allowbreak{}thesis\_rrf\_example.json}). Khi vector tên điều yếu, RRF có thể đẩy xuống một kết quả mà BM25 xếp đầu.

\section{Tác tử tự kiểm tra quy định hiện hành}
\label{section:5.2}

\subsection{Vấn đề}
Bộ Pháp điển là nguồn chính thống nhưng được cập nhật theo chu kỳ, nên có thể chưa phản ánh văn bản mới ban hành. Ví dụ, các mức phạt vi phạm giao thông trong Nghị định 100/2019/NĐ-CP đã được thay thế bởi Nghị định 168/2024/NĐ-CP. Trong thiết kế ban đầu, tác tử tự quyết định có gọi công cụ tìm kiếm web hay không. Thực tế cho thấy khi đã thấy một điều luật trong Pháp Điển có vẻ khớp với câu hỏi, mô hình hiếm khi gọi thêm tìm kiếm web, dẫn tới trả lời theo mức phạt cũ một cách tự tin.

Vấn đề thứ hai là cách trích dẫn. Mỗi điều trong Pháp Điển có mã sắp xếp dạng ``39.13.NĐ.75.8''; mô hình thường đọc nhầm phần ``NĐ.75'' thành ``Nghị định 75'' hoặc ``39.13'' thành ``Điều 39.13'', tạo ra trích dẫn không tồn tại.

\subsection{Giải pháp}
Đồ án bổ sung một luật buộc kiểm tra web nằm ngoài quyết định của mô hình. Sau mỗi vòng gọi công cụ, trước lượt gọi mô hình tiếp theo, nếu tác tử chưa gọi tìm kiếm web thì hệ thống kiểm tra ba điều kiện: (i) tác tử đã tra Pháp Điển nhưng không có kết quả; (ii) câu hỏi chứa các dấu hiệu dễ thay đổi theo văn bản mới như ``phạt'', ``bao nhiêu tiền'', ``lệ phí'', ``thuế'', ``mức lương'', ``trợ cấp'', ``mức hưởng'', ``hiện nay'', ``hiện hành'', ``còn hiệu lực'' hoặc ``năm 20xx''; hoặc (iii) điểm liên quan cao nhất của kết quả Pháp Điển dưới 0,5. Nếu một điều kiện thoả mãn, hệ thống chèn một chỉ dẫn nêu lý do và buộc lượt gọi mô hình đó phải gọi \texttt{web\_search}, mỗi câu hỏi tối đa một lần. Việc kiểm tra được đặt trước khi mô hình viết câu trả lời, không phải sau, vì câu trả lời được truyền tới người dùng ngay khi mô hình sinh; nếu kiểm tra sau thì bản nháp đã hiển thị sẽ phải bị thay thế. Kết quả web được lọc theo độ tin cậy tên miền: các trang chính thống và thư viện pháp luật lớn được xếp mức cao, các trang blog, mạng xã hội, chia sẻ tài liệu bị loại. Lời nhắc hệ thống quy định: Pháp Điển là nguồn chính về nội dung, nhưng nếu nguồn web tin cậy cao cho thấy văn bản đã bị thay thế hoặc mức tiền đã đổi thì phải nêu quy định hiện hành kèm nguồn và ghi chú quy định cũ.

Để giải quyết vấn đề trích dẫn, khi đưa điều luật cho mô hình, hệ thống tách mã sắp xếp khỏi tên điều thành trường riêng \texttt{ma\_phap\_dien}, và lời nhắc yêu cầu số điều cùng số hiệu văn bản phải lấy nguyên văn từ trường \texttt{vbqppl} (ví dụ ``Điều 8 Nghị định số 100/2019/NĐ-CP''). Ngoài ra, mỗi điều luật được gắn chú thích hiệu lực tra từ sổ cái văn bản (phần \ref{section:5.5}); điều thuộc văn bản đã hết hiệu lực bị nhân điểm với hệ số 0,4 và được ghi chú để mô hình không dùng làm căn cứ chính.

\subsection{Kết quả đạt được}
Hành vi gọi công cụ được đánh giá trên 23 tình huống có kỳ vọng ghi sẵn, chia thành bảy nhóm: định nghĩa, phạm vi áp dụng, mức phạt, câu hỏi ngoài dữ liệu, quan hệ giữa điều luật, đọc toàn văn và trò chuyện. Hai mô hình tác tử được chạy trên cùng mã nguồn; kết quả chi tiết ở Bảng \ref{tab:pl_tool} (Phụ lục \ref{appendix:A}). Với GPT-4.1, cấu hình mặc định của hệ thống, 21 trên 23 tình huống đạt mọi kỳ vọng; tỷ lệ gọi đủ công cụ bắt buộc là 97,6\%, công cụ đầu tiên đúng kỳ vọng và luật buộc web kích hoạt đúng ở mọi tình huống, không có lượt gọi công cụ bị cấm hay mã điều bịa đặt, và một tình huống chạm trần vòng lặp. Hai tình huống chưa đạt gồm một câu trả lời thiếu từ khoá kỳ vọng và một câu hỏi ngoài dữ liệu mà tác tử không gọi tìm kiếm web: luật buộc web không kích hoạt vì câu hỏi không chứa dấu hiệu nào trong danh sách và Pháp Điển vẫn trả về điều có điểm từ 0,5 trở lên. Với GPT-4o-mini, 19 trên 23 tình huống đạt; mô hình này bỏ lỡ tìm kiếm web ở cả ba câu hỏi ngoài dữ liệu và truyền mã điều bịa đặt ở 8,3\% lượt gọi đồ thị hoặc đọc toàn văn, nhưng chi phí cho cả bộ chỉ 0,09 USD so với 1,62 USD của GPT-4.1. Kết quả cho thấy luật buộc web hoạt động đúng khi điều kiện được thoả, còn các câu hỏi ngoài dữ liệu không mang dấu hiệu thời gian vẫn phụ thuộc vào khả năng tự nhận biết của mô hình.

Bộ đánh giá chạy mỗi tình huống một lần nên chưa đo được độ dao động giữa các lần chạy; với 23 tình huống, chênh lệch hai tình huống giữa hai mô hình chỉ mang tính tham khảo.

\section{Đồ thị pháp luật: dẫn chiếu và kiểm tra ngoại lệ}
\label{section:5.graph}

\subsection{Vấn đề}
Một điều luật thường không tự đủ nghĩa. Điều 8 Luật Hôn nhân và gia đình 2014 về điều kiện kết hôn dẫn chiếu tới các điểm a, b, c, d khoản 2 Điều 5 của luật này, nơi quy định các trường hợp cấm kết hôn. Tương tự, Điều 125 Bộ luật Lao động 2019 nêu các trường hợp được áp dụng hình thức sa thải, còn các trường hợp không được xử lý kỷ luật lao động nằm ở khoản 4 Điều 122. Tìm kiếm theo câu hỏi thường trả về điều chính, vì câu hỏi dùng từ ngữ của điều chính, còn điều chứa ngoại lệ có từ ngữ khác. Nếu tác tử chỉ đọc điều chính, câu trả lời có thể kết luận ``được phép'' trong khi tình huống của người hỏi thuộc trường hợp bị cấm.

Ở phiên bản trước, đồ thị tri thức chỉ có cạnh dẫn chiếu tạo từ mã Pháp Điển xuất hiện trong nội dung, nên tạo được rất ít cạnh, và công cụ mở rộng đồ thị chỉ được gọi khi mô hình tự quyết định.

\subsection{Giải pháp}
Giải pháp gồm ba phần. Thứ nhất, bộ phân tích dẫn chiếu (phần \ref{subsection:4.data.graph}) tạo 24.718 cạnh \texttt{THAM\_CHIEU} từ các dẫn chiếu dạng ``khoản 2 Điều 5 của Luật này'', trong đó 1.250 cạnh nằm trong câu mang nghĩa ngoại lệ. Mỗi khoản được gắn cờ ngoại lệ theo định nghĩa ở phần \ref{subsection:4.data.graph}, tạo thành 628 khoản ngoại lệ; 1.104 điều được đánh dấu là điều ngoại lệ.

Thứ hai, công cụ \texttt{vector\_search} tự bổ sung hai loại nội dung sau khi xếp hạng lại, không cần mô hình gọi thêm công cụ. \emph{Kiểm tra ngoại lệ} lấy năm điều đầu danh sách, tìm các điều nối với chúng qua cạnh dẫn chiếu theo cả hai chiều hoặc cạnh liên quan mà là điều ngoại lệ hoặc được dẫn chiếu trong ngữ cảnh ngoại lệ, rồi đưa tối đa ba điều vào kết quả. Mỗi điều được bổ sung chỉ gồm các khoản ngoại lệ và các khoản được dẫn chiếu, kèm trường \texttt{ngoai\_le\_cua} ghi điều chính mà nó giới hạn và lời nhắc đối chiếu trước khi kết luận; ưu tiên dẫn chiếu trong ngữ cảnh ngoại lệ, sau đó tới dẫn chiếu thường và cạnh liên quan. \emph{Kèm điều được dẫn chiếu} lấy tối đa bốn điều hoặc khoản mà ba điều đầu dẫn chiếu tới, mỗi mục tối đa 1.500 ký tự, kèm trường \texttt{duoc\_dan\_chieu\_tu}. Cả hai chỉ thêm điều chưa có trong danh sách. Hệ thống không tự kết luận ngoại lệ có áp dụng hay không; việc đối chiếu với tình huống do tác tử thực hiện theo lời nhắc hệ thống.

Thứ ba, khi một điều dài hơn 4.000 ký tự được rút gọn trước khi đưa cho mô hình, các khoản ngoại lệ luôn được giữ lại cùng các khoản khớp câu hỏi nhất. Công cụ \texttt{graph\_traverse} vẫn được giữ cho trường hợp mô hình cần mở rộng thêm, với trọng số theo loại quan hệ: dẫn chiếu đi ra 0,8, được dẫn chiếu 0,65, liên quan 0,5 và bước dẫn chiếu thứ hai 0,4.

\subsection{Kết quả đạt được}
Hai ví dụ sau lấy từ kết quả thật của công cụ \texttt{vector\_search} với 10 kết quả, chạy bằng \texttt{scripts/\allowbreak{}thesis\_graph\_examples.py} và lưu tại \texttt{evaluation/\allowbreak{}results/\allowbreak{}thesis\_graph\_examples.json}.

Với truy vấn ``điều kiện kết hôn'', Điều 8 Luật số 52/2014/QH13 đứng đầu. Điều 5 của luật này (Bảo vệ chế độ hôn nhân và gia đình) không có trong 10 kết quả tìm kiếm nhưng được bổ sung ở vị trí 11: khi tắt kiểm tra ngoại lệ, điều này vẫn được kèm theo dưới dạng điều được dẫn chiếu từ Điều 8 (``dẫn chiếu điểm a, b, c, d khoản 2''); khi bật, nó được gắn nhãn ngoại lệ của Điều 8 vì cạnh dẫn chiếu này nằm trong ngữ cảnh ngoại lệ. Nội dung trả cho tác tử gồm khoản 2, trong đó điểm d cấm ``kết hôn hoặc chung sống như vợ chồng giữa những người cùng dòng máu về trực hệ; giữa những người có họ trong phạm vi ba đời\ldots''. Với tình huống người hỏi muốn kết hôn với con gái của chú ruột, đây là căn cứ quyết định mà Điều 8 không chứa.

Với truy vấn ``tự ý bỏ việc 5 ngày bị sa thải'', Điều 125 Bộ luật số 45/2019/QH14 đứng đầu. Khi bật kiểm tra ngoại lệ, Điều 122 (Nguyên tắc, trình tự, thủ tục xử lý kỷ luật lao động) được bổ sung với nhãn ngoại lệ của Điều 125; khi tắt, điều này không có trong kết quả. Trong đồ thị, cạnh dẫn chiếu giữa hai điều đi từ Điều 122 tới Điều 125 (điểm c khoản 4 Điều 122 dẫn chiếu khoản 1 và khoản 2 Điều 125), nên phép kèm điều được dẫn chiếu, vốn chỉ đi theo chiều ra từ điều đứng đầu, không tìm được; kiểm tra ngoại lệ xét cả hai chiều và Điều 122 có khoản ngoại lệ nên được chọn. Nội dung trả cho tác tử gồm khoản 4 với điểm d ``Người lao động nữ mang thai; người lao động nghỉ thai sản, nuôi con dưới 12 tháng tuổi'', là căn cứ cho tình huống người sử dụng lao động muốn sa thải lao động nữ đang mang thai.

Để đo ở quy mô lớn hơn, bộ đánh giá \texttt{eval\_exceptions.py} lấy từ đồ thị các cặp điều chính A và điều ngoại lệ B nối nhau qua cạnh dẫn chiếu hoặc liên quan. Claude Sonnet đọc cả hai điều, chỉ giữ các cặp mà B thực sự giới hạn điều A cho phép, và viết một câu hỏi đời thường mà A trả lời nhưng không nhắc tới điều kiện ngoại lệ. Câu hỏi được viết lại như trong hệ thống rồi đưa vào \texttt{vector\_search}, so sánh khi tắt và bật kiểm tra ngoại lệ. Bảng \ref{tab:ngoaile} trình bày kết quả trên 30 cặp.

\input{tables/tab_ngoaile}

Khi bật kiểm tra ngoại lệ, điều ngoại lệ có trong kết quả ở 14 trên 30 trường hợp so với 12 khi tắt; xét riêng các trường hợp điều chính có mặt, tỷ lệ là 9/15 so với 8/15. Kiểm tra ngoại lệ không làm mất điều chính và chỉ thêm trung bình 0,4 kết quả. Tuy nhiên, điều chính chỉ được tìm thấy ở 50\% số trường hợp, và kiểm tra ngoại lệ chỉ có tác dụng khi điều chính đã có trong danh sách, nên chất lượng truy xuất đang là giới hạn chính. Với 30 cặp, chênh lệch hai trường hợp chưa đủ để kết luận mức cải thiện; bộ đánh giá cũng chưa đo việc tác tử có dùng đúng ngoại lệ trong câu trả lời hay không.

\section{Kiểm chứng căn cứ trong rà soát văn bản}
\label{section:5.3}

\subsection{Vấn đề}
Rà soát hợp đồng là tác vụ có rủi ro cao hơn hỏi đáp: hệ thống phải tự tìm ra điều khoản có vấn đề trong một văn bản dài và đưa ra kết luận ``trái luật'', vốn có thể ảnh hưởng tới quyết định ký kết của người dùng. Nếu để LLM tự kết luận và tự nêu căn cứ, mô hình có thể nêu một điều luật đúng về nội dung nhưng sai số điều, hoặc nêu điều luật hoàn toàn không tồn tại. Trong quá trình phát triển, đồ án ghi nhận ba lỗi cụ thể: (i) phần giải thích dẫn ``Điều 98 Bộ luật Lao động'' trong khi căn cứ được chọn lại là Điều 91 về lương tối thiểu; (ii) mô hình nhỏ chép sai một chữ số trong mã điều luật dài từ 40 đến 120 chữ số, khiến hệ thống không khớp được căn cứ và bỏ toàn bộ cảnh báo; (iii) khi có cả nguồn web và Pháp Điển cùng nói về một điều luật, mô hình chỉ dẫn bài báo trên web.

\subsection{Giải pháp}
Đồ án thiết kế luồng rà soát gồm hai lần gọi LLM tách bạch, xen giữa là bước tra cứu do mã nguồn kiểm soát, và một bước kiểm chứng hoàn toàn bằng mã nguồn sau lần gọi thứ hai.

Lần gọi thứ nhất yêu cầu mô hình đọc văn bản (được bọc trong thẻ riêng và được chỉ dẫn là dữ liệu, bỏ qua mọi câu lệnh bên trong để chống chèn lệnh) và trả về JSON gồm tóm tắt, các bên, ý chính, các điều khoản rủi ro. Mỗi điều khoản rủi ro kèm hai truy vấn tra cứu: một cụm từ pháp lý và tiêu đề điều luật dự kiến điều chỉnh. Mã nguồn chạy cả hai truy vấn trên tìm kiếm lai (phần \ref{section:5.1}), gộp theo mã điều, và tra thêm web theo đúng luật của tác tử hỏi đáp (phần \ref{section:5.2}). Việc tra hai truy vấn giúp tìm được điều luật mà một cụm từ duy nhất bỏ sót; trong thử nghiệm, Điều 98 Bộ luật Lao động về tiền lương làm thêm giờ chỉ được tìm thấy nhờ truy vấn theo tiêu đề điều.

Lần gọi thứ hai đưa cho mô hình từng điều khoản cùng danh sách nguồn ứng viên của riêng điều khoản đó. Mỗi nguồn được gán mã ngắn \texttt{N1}, \texttt{N2}\ldots{} thay cho mã điều luật dài, khắc phục lỗi (ii). Mô hình phải trả về kết luận, các căn cứ theo mã ngắn kèm đoạn trích nguyên văn, và lời giải thích dẫn căn cứ bằng số thứ tự trong ngoặc vuông, không được viết số điều trong lời giải thích.

Bước kiểm chứng sau đó áp dụng sáu quy tắc. Thứ nhất, căn cứ chỉ được chấp nhận nếu mã của nó thuộc danh sách ứng viên của chính điều khoản đó. Thứ hai, đoạn trích phải có thật trong nội dung nguồn; phép so khớp chuẩn hoá chữ hoa thường, dấu ngoặc và khoảng trắng, cho phép dấu ``\ldots'' bỏ đoạn, và yêu cầu mỗi đoạn dài từ 15 ký tự trở lên khớp liên tục ít nhất 85\% độ dài; đoạn trích không khớp bị bỏ nhưng căn cứ vẫn được giữ. Thứ ba, nếu nguồn web nêu một số điều mà Pháp Điển có trong danh sách ứng viên, điều Pháp Điển đó được tự động chèn lên trước nguồn web và các số tham chiếu được đánh lại, khắc phục lỗi (iii). Thứ tư, kết luận ``trái luật'' phải có ít nhất một căn cứ Pháp Điển hoặc nguồn web độ tin cậy cao, nếu không bị hạ xuống ``rủi ro''. Thứ năm, điều luật được nhắc trong lời giải thích mà không nằm trong căn cứ được liệt kê riêng để giao diện ghi chú ``chưa được kiểm chứng'', khắc phục lỗi (i). Thứ sáu, điều khoản không còn căn cứ hợp lệ nào không được hiển thị như cảnh báo mà được chuyển sang mục ``Điều khoản nên xem lại''.

Trên giao diện, các điều luật đã tra được trả về theo đúng định dạng của tính năng hỏi đáp, nên khung căn cứ pháp lý được dùng lại; số tham chiếu trong lời giải thích và thẻ căn cứ đều mở đúng điều luật. Toàn bộ quy tắc trên được bao phủ bởi 36 kiểm thử đơn vị (Bảng \ref{tab:tc_rasoat}).

\subsection{Kết quả đạt được}
Các quy tắc kiểm chứng được kiểm tra bằng 36 kiểm thử đơn vị; Bảng \ref{tab:tc_rasoat} liệt kê các trường hợp ứng với ba lỗi đã ghi nhận ở trên (trường hợp 1, 6, 7 và 8) và với quy tắc hạ mức kết luận (trường hợp 3 đến 5). Đồ án đã chạy thử luồng rà soát trên các hợp đồng mẫu do người thực hiện soạn, có cài sẵn điều khoản trái luật, nhưng không lưu lại kết quả của các lần chạy này, nên không báo cáo số lượng cảnh báo hay tỷ lệ phát hiện.

Thử nghiệm cũng cho thấy hai giới hạn. Kết quả có độ dao động giữa các lần chạy vì cụm từ tra cứu do mô hình tự đặt; khi một lần chạy không tìm được căn cứ, hệ thống chuyển điều khoản sang mục ``nên xem lại'' thay vì bịa căn cứ, đổi lại có thể bỏ sót cảnh báo đúng. Ngoài ra, hệ thống bảo đảm căn cứ là điều luật có thật và đoạn trích là nguyên văn, nhưng không bảo đảm việc áp dụng điều luật vào tình huống luôn đúng về mặt pháp lý, ví dụ áp dụng trần lãi suất của Bộ luật Dân sự cho lãi chậm trả trong hợp đồng thương mại là vấn đề còn nhiều ý kiến. Vì vậy mọi kết quả đều kèm lưu ý không thay thế ý kiến của luật sư.

\section{Bộ nhớ dài hạn bảo vệ quyền riêng tư}
\label{section:5.4}

\subsection{Vấn đề}
Câu trả lời pháp lý thường phụ thuộc vào hoàn cảnh cá nhân như tuổi, nghề nghiệp, loại hợp đồng. Nếu hệ thống không nhớ, người dùng phải nhắc lại thông tin ở mỗi cuộc hội thoại; nếu nhớ không kiểm soát, hệ thống có thể lưu những thông tin rất nhạy cảm như số căn cước, số tài khoản, tiền án, tình trạng sức khoẻ, vốn hay xuất hiện trong câu hỏi pháp luật. Ngoài ra, thông tin thay đổi theo thời gian: người dùng chuyển trường, đổi việc, nên bộ nhớ phải cập nhật thay vì tích luỹ các phiên bản mâu thuẫn.

\subsection{Giải pháp}
Đồ án áp dụng kiến trúc hai giai đoạn của Mem0 \cite{chhikara2025mem0}: sau mỗi lượt hội thoại, một mô hình nhỏ trích xuất các thông tin bền vững; với mỗi thông tin, hệ thống tìm tối đa năm ký ức tương đồng của người dùng và để mô hình quyết định thêm, cập nhật, xoá hay bỏ qua. Toàn bộ quá trình chạy nền sau khi đã trả lời. Đồ án bổ sung bốn điểm so với kiến trúc gốc.

Thứ nhất, ký ức được phân thành bốn loại với độ phai khác nhau khi truy xuất: thông tin cốt lõi (năm sinh, nghề nghiệp) và thói quen trình bày không phai; sự kiện (bị nợ lương tháng này) có chu kỳ bán rã 30 ngày; kiến thức về người dùng có chu kỳ bán rã 120 ngày. Điểm truy xuất được tính bằng
\begin{equation}\label{eq:memscore}
    \mathrm{score} = \mathrm{sim}\cdot\left((1-w) + w\cdot 0{,}5^{\,\Delta t / T_{1/2}}\right)
\end{equation}
với $\mathrm{sim}$ là độ tương đồng cosine, $\Delta t$ là tuổi của ký ức tính theo ngày, $T_{1/2}$ là chu kỳ bán rã theo loại và $w = 0{,}4$.

Thứ hai, quyền riêng tư được bảo vệ ở hai lớp. Mô hình trích xuất gắn cờ nhạy cảm cho từng thông tin; một lớp biểu thức chính quy làm lưới dự phòng, loại các dãy số dạng số căn cước, số điện thoại, số tài khoản, số thẻ và nhận diện các chủ đề nhạy cảm như tố tụng hình sự, sức khoẻ, xu hướng tính dục, tôn giáo, chính trị. Người dùng chọn một trong ba chế độ: không lưu thông tin nhạy cảm (mặc định), lưu đầy đủ, hoặc tắt bộ nhớ. Người dùng xem được mọi ký ức, xoá từng mục hoặc xoá toàn bộ.

Thứ ba, mọi thao tác trên bộ nhớ đều kiểm tra quyền: người dùng đã đăng nhập chỉ thao tác được bộ nhớ của chính mình, mã người dùng gửi lên mà không có token hợp lệ không được tin, trừ mã ẩn danh do chính trình duyệt sinh ra. Thứ tư, bộ nhớ và vector của nó được lưu cùng một tài liệu MongoDB thay vì tách ra Milvus như thiết kế ban đầu; việc tách hai kho đã gây lệch dữ liệu khi Milvus được khởi tạo lại.

\subsection{Kết quả đạt được}
Bộ nhớ được đánh giá với mô hình, MongoDB thật trên 14 kịch bản hội thoại gồm bốn nhóm: trích xuất (6 kịch bản, như thông tin cơ bản, trò chuyện không có gì để lưu, câu hỏi pháp luật chung không phải thông tin cá nhân), cập nhật (4 kịch bản, như đổi trường, đổi việc, phủ định cần xoá, diễn đạt lại cần bỏ qua), nhạy cảm (2 kịch bản với số căn cước và số tài khoản) và truy xuất (2 kịch bản). Bộ đánh giá được chạy bốn lần (Bảng \ref{tab:pl_mem_runs}, Phụ lục \ref{appendix:A}). Ở cả bốn lần, 14 kịch bản đều đạt, không có lần lộ dữ liệu nhạy cảm nào và không có ký ức trùng lặp hay lỗi thời. Ba lần chạy gần nhất có cùng phân bố thao tác gồm 18 lần thêm, 1 lần cập nhật, 2 lần xoá và 1 lần bỏ qua; lần đầu có 17 lần thêm. Trung vị thời gian xử lý mỗi kịch bản từ 2,50 đến 4,93 giây. Đồ án báo trung vị thay cho trung bình vì một lần chạy có một kịch bản kéo dài 878 giây; mã nguồn bộ nhớ không thay đổi giữa các lần chạy và hai lần chạy lại không lặp lại hiện tượng này.

\section{Cập nhật văn bản pháp luật và theo dõi hiệu lực}
\label{section:5.5}

\subsection{Vấn đề}
Dữ liệu Pháp Điển chỉ được tải một lần lúc xây dựng hệ thống, trong khi văn bản pháp luật mới được ban hành liên tục. Nếu không có cách cập nhật, hệ thống sẽ ngày càng lạc hậu; nếu cập nhật tự động không qua kiểm tra, dữ liệu sai do lỗi trích xuất PDF sẽ đi thẳng vào câu trả lời. Ngoài ra, hệ thống cần biết văn bản nào đã hết hiệu lực, đã bị sửa đổi hay thay thế để cảnh báo người dùng.

\subsection{Giải pháp}
Đồ án xây dựng mô-đun quản trị dữ liệu văn bản với quy trình qua khu chờ duyệt. Quản trị viên thêm tối đa 30 văn bản mỗi lần từ trang văn bản trên Công báo, đường dẫn PDF hoặc file PDF. Việc tải từ đường dẫn tuân thủ \texttt{robots.txt}, giãn cách giữa các lần tải và chặn địa chỉ nội bộ. Bộ trích xuất PDF lọc chữ có cỡ nhỏ hơn 80\% cỡ chữ nội dung để loại chú thích cuối trang, tránh lỗi số chú thích ``1'' dính vào ``Điều 5'' thành ``Điều 51'' trong văn bản hợp nhất. Bộ phân tích cấu trúc nhận diện số hiệu, loại văn bản, ngày ban hành, ngày hiệu lực và các cấp phần, chương, mục, điều theo đúng văn bản đó; nội dung nằm trong ngoặc kép, tức là phần văn bản sửa đổi trích lại điều cũ, không bị coi là cấu trúc của văn bản đang đọc. Bộ phân tích cũng nhận diện các câu nêu quan hệ sửa đổi, thay thế với văn bản khác.

Mỗi văn bản được so sánh với dữ liệu hiện có để đánh dấu điều mới, điều thay đổi và điều giữ nguyên, rồi chờ quản trị viên kiểm tra, sửa và duyệt. Dữ liệu chỉ được ghi vào kho tra cứu sau khi duyệt; khu chờ duyệt không ảnh hưởng tới câu trả lời. Khi ghi, nội dung cũ của điều bị thay đổi được lưu vào bảng lịch sử; điều thuộc văn bản ngoài Pháp Điển được gán mã tạm và khi Pháp Điển cập nhật điều tương ứng, mã tạm được ánh xạ sang mã chính thức để các trích dẫn cũ vẫn mở được. Sổ cái văn bản lưu trạng thái hiệu lực ở mức văn bản và mức điều; khi truy xuất, mỗi điều luật được gắn chú thích hiệu lực và điều thuộc văn bản hết hiệu lực bị hạ điểm như mô tả ở phần \ref{section:5.2}.

\subsection{Kết quả đạt được}
Quy trình cập nhật đã hoạt động đầy đủ trên giao diện quản trị gồm bốn thẻ: thêm văn bản, chờ duyệt, sổ cái văn bản và cập nhật Pháp Điển. Chức năng quản trị chỉ mở cho tài khoản có email đã xác minh và được cấp vai trò trong MongoDB. Mô hình hiện tại lưu một trạng thái hiệu lực cho mỗi điều; việc lưu nội dung điều theo nhiều phiên bản gắn với khoảng thời gian hiệu lực, để xử lý trọn vẹn các trường hợp sửa đổi, thay thế và hợp nhất, được trình bày như hướng phát triển ở Chương \ref{chapter:conclusion}.

\section{Trực quan hoá hoạt động của tác tử}
\label{section:5.6}

\subsection{Vấn đề}
Một lượt trả lời của tác tử mất từ vài giây tới hơn mười giây. Trong khoảng thời gian đó, giao diện chat thông thường chỉ hiển thị biểu tượng chờ, người dùng không biết hệ thống đang làm gì và dễ nghĩ hệ thống bị treo. Đồng thời, kiến trúc nhiều thành phần (tiền xử lý, bộ đệm, bộ nhớ, bốn công cụ) khó giải thích cho người không chuyên.

\subsection{Giải pháp}
Máy chủ phát các sự kiện trạng thái có cấu trúc trong luồng SSE, mỗi sự kiện có trường \texttt{stage} cho biết bước hiện tại (tiếp nhận, chuẩn hoá, bộ nhớ, bộ đệm, gọi công cụ, tổng hợp) và tên công cụ, số vòng lặp. Trang hỏi đáp hiển thị các sự kiện này dưới dạng dòng tiến trình. Trang văn phòng luật 3D ánh xạ mỗi thành phần thành một nhân vật như mô tả ở phần \ref{subsection:4.3.3}: khi máy chủ báo đang gọi \texttt{vector\_search}, nhân vật thủ thư điều luật làm việc tại thư viện Pháp Điển; khi gọi \texttt{graph\_traverse}, chuyên viên liên kết làm việc tại bảng liên kết; khi bộ đệm trúng, ``sổ tay đã trả lời'' được sử dụng thay cho cả quy trình tra cứu. Trang hỗ trợ ba góc nhìn (toàn cảnh, nhập vai, mắt nhân vật) và một chế độ dữ liệu mẫu chạy không cần máy chủ để trình diễn.

\subsection{Kết quả đạt được}
Cùng một luồng sự kiện phục vụ cả giao diện chat và văn phòng 3D mà không cần điểm cuối riêng, vì các sự kiện vẫn giữ trường \texttt{text} cho giao diện cũ. Văn phòng 3D biến luồng xử lý trừu tượng của tác tử thành hình ảnh trực quan, giúp người dùng và người đánh giá thấy được tác tử đang tra Pháp Điển, đi theo đồ thị hay tìm kiếm web ở từng thời điểm.

Chương này đã trình bày bảy giải pháp chính của đồ án. Tìm kiếm lai với bộ tách từ cặp âm tiết, RRF có trọng số và xếp hạng lại theo độ phủ từ có nghĩa nâng P@10 từ 0,368 lên 0,544 và MRR từ 0,280 lên 0,471 so với cấu hình ban đầu, dù chưa tốt hơn BM25 đơn lẻ một cách có ý nghĩa thống kê. Luật buộc kiểm tra web và cách tách mã sắp xếp khỏi trích dẫn giúp tác tử GPT-4.1 gọi đủ công cụ bắt buộc ở 97,6\% trường hợp. Đồ thị dẫn chiếu và kiểm tra ngoại lệ bổ sung các khoản cấm, ngoại lệ mà tìm kiếm bỏ sót, với giới hạn là điều chính phải được tìm thấy trước. Cơ chế kiểm chứng căn cứ bảo đảm mọi cảnh báo trong rà soát văn bản dựa trên điều luật có trong danh sách đã tra và đoạn trích nguyên văn. Bộ nhớ dài hạn đạt toàn bộ kịch bản kiểm thử mà không lộ dữ liệu nhạy cảm. Quy trình qua khu chờ duyệt cho phép cập nhật văn bản mới một cách có kiểm soát, và văn phòng 3D trực quan hoá hoạt động của tác tử. Chương tiếp theo tổng kết đồ án và đề xuất hướng phát triển.

\end{document}
